The displayed bands are posterior distributions of the draw-wise aggregate criteria under each model's native conditional-quantile construction. Q--DESN and exQ--DESN preserve the fitted posterior readout-draw identity across rolling origins and forecast leads. DQLM and exDQLM combine the origin--lead latent quantile marginals produced by the rolling state-update interface using the pre-specified product coupling. The resulting intervals are valid for these stated model-specific contracts; they do not impose a common cross-origin copula, and their widths should therefore be interpreted as native posterior uncertainty rather than as a model-free measure of repeated-simulation variability.

A paired coupling sensitivity analysis preserved every finite origin--lead marginal distribution and left the posterior metric means unchanged to numerical precision. Interval widths were sensitive to the imposed dependence in the representative cells, but no tested winner--runner-up interval-overlap conclusion changed. The native intervals are consequently retained, with the coupling analysis serving as sensitivity evidence rather than as a replacement estimator.
